\documentclass[aspectratio=169]{beamer}

\usepackage{ctex}
\usepackage{amsmath, amsfonts, amssymb, bm}
\usepackage{graphicx}
\usepackage{hyperref}
\usepackage{booktabs}
\usepackage{array}
\usepackage{tabularx}
\usepackage{ragged2e}

\usetheme{Berlin}
\usecolortheme{whale}
\useinnertheme[shadow=true]{rounded}
\usefonttheme[onlymath]{serif}
\setbeamercovered{transparent}
\setbeamertemplate{navigation symbols}{}
\setbeamertemplate{bibliography item}[text]
\setbeamercolor{bibliography item}{fg=black}
\setbeamercolor{bibliography entry author}{fg=black}
\setbeamercolor{bibliography entry title}{fg=black}
\setbeamercolor{bibliography entry location}{fg=black}
\setbeamercolor{bibliography entry note}{fg=black}
\hypersetup{hidelinks}

\newcolumntype{Y}{>{\RaggedRight\arraybackslash}X}
\newcommand{\good}[1]{\textcolor{green!45!black}{#1}}
\newcommand{\warn}[1]{\textcolor{orange!80!black}{#1}}
\newcommand{\bad}[1]{\textcolor{red!70!black}{#1}}

\setbeamertemplate{footline}{%
 \leavevmode%
 \hbox{%
 \begin{beamercolorbox}[wd=.31\paperwidth,ht=2.25ex,dp=1ex,center]{author in head/foot}
 \insertshortauthor
 \end{beamercolorbox}%
 \begin{beamercolorbox}[wd=.38\paperwidth,ht=2.25ex,dp=1ex,center]{title in head/foot}
 \insertshorttitle
 \end{beamercolorbox}%
 \begin{beamercolorbox}[wd=.31\paperwidth,ht=2.25ex,dp=1ex,center]{date in head/foot}
 \insertshortdate \quad \insertframenumber/\inserttotalframenumber
 \end{beamercolorbox}%
 }
}

\title[晶体带隙预测阶段性汇报]{晶体带隙预测项目阶段性汇报}
\subtitle{从模型增强到数据质量诊断：HZW 与 Materials Project 复现分析}
\author[刘嘉瑞]{汇报人：刘嘉瑞}
\institute[GNN Band Gap]{任务：晶体 band gap 回归预测\\模型：CGCNN / ALIGNN / M3GNet / Virtual Node / Attentive Dual Slot}
\date{2026年5月}

\begin{document}

\begin{frame}
 \titlepage
\end{frame}

\begin{frame}{目录}
 \tableofcontents
\end{frame}

\section{背景与主线}

\begin{frame}{汇报主线}
\small
\begin{block}{研究问题}
局部晶体 GNN 是否缺少全局信息？在 ALIGNN 上加入 virtual node、attention slot 和物理约束，能否改善 band gap 预测？
\end{block}

\begin{block}{实验路径}
先在鸿之微 HZW 数据集上验证 E1--E13，再迁移到 Materials Project full 和 MP-HZW matched 数据。
\end{block}

\begin{alertblock}{当前核心结论}
模型排序不能脱离标签源讨论；MP summary band gap 与 HZW band gap 的冲突足以改变模型结论。
\end{alertblock}
\end{frame}

\begin{frame}{任务背景：晶体 band gap 预测}
\small
\begin{itemize}
  \item 输入：晶体结构，包括原子种类、坐标、晶格和邻接关系。
  \item 输出：材料带隙 \texttt{band\_gap}，单位 eV。
  \item 经典晶体 GNN 主要依赖局部消息传递：邻居原子、键长、键角和局部配位环境。
  \item 本项目假设：band gap 虽强依赖局部化学，但晶胞尺度、长程结构和全局汇聚也可能影响预测。
  \item 因此主线是：忠实复现基线后，逐步加入超节点、注意力槽位和物理一致性约束。
\end{itemize}
\end{frame}

\begin{frame}{三个基线模型}
\scriptsize
\centering
\setlength{\tabcolsep}{4pt}
\begin{tabularx}{\textwidth}{p{2.1cm} Y Y}
\toprule
模型 & 核心思想 & 本项目中的作用 \\
\midrule
CGCNN & 将晶体表示为原子图，边特征主要来自邻居距离。 & 最基础晶体图卷积 baseline，简单但几何表达较弱。 \\
ALIGNN & 在原子图之外引入 line graph，显式建模键角和二阶几何关系。 & 本项目最强、最稳定的 baseline，也是后续增强模型主干。 \\
M3GNet & 强调三体几何和通用材料图表示。 & 材料建模能力强，但在当前 band gap 任务中没有超过 ALIGNN。 \\
\bottomrule
\end{tabularx}
\end{frame}

\begin{frame}{数据与阶段路线}
\scriptsize
\centering
\setlength{\tabcolsep}{4pt}
\begin{tabularx}{\textwidth}{p{2.0cm} p{2.5cm} Y Y}
\toprule
阶段 & 数据范围 & 目的 & 关键发现 \\
\midrule
HZW & 95,497 条，固定 train/val/test。 & 验证模型增强是否有效。 & E4 稳定提升，E12/E13 进一步超过 E4。 \\
MP full & 154,875 条 MP bandgap。 & 检验公开数据泛化。 & E4 重新最稳，E7/E8/E12 未超过 E4。 \\
MP-HZW matched & 92,826 条匹配样本。 & 一比一对齐 HZW 来源。 & 发现 MP summary 与 HZW band gap 标签冲突。 \\
Clean matched & 84,830 条类别一致样本。 & 排除金属/非金属冲突后重训。 & E4 仍最稳，E12/E1 待最终同步。 \\
\bottomrule
\end{tabularx}
\end{frame}

\section{方法}

\begin{frame}{实验矩阵：E1--E13}
\tiny
\centering
\setlength{\tabcolsep}{3pt}
\begin{tabularx}{\textwidth}{p{0.7cm} p{3.1cm} Y p{2.5cm}}
\toprule
编号 & 模型/方案 & 核心作用 & 当前判断 \\
\midrule
E1 & CGCNN & 基础晶体图卷积 baseline。 & 弱于 ALIGNN \\
E2 & ALIGNN & 原子图 + line graph / 键角信息。 & 最强 baseline \\
E3 & M3GNet & 三体几何和通用材料图。 & 本任务未超过 ALIGNN \\
E4 & ALIGNN + 朴素超节点 & 全局汇聚。 & 最稳定增强 \\
E5 & ALIGNN + Transformer & 全局注意力 only。 & 整体负优化/弱收益 \\
E6 & ALIGNN + 物理约束 & CBM/VBM 辅助监督。 & HZW 有收益 \\
E7 & 1 个 attentive 超节点 & 单槽位注意力超节点。 & 不稳定，HZW 负优化 \\
E8 & 2 个 attentive 超节点 & 双槽位 + slot attention。 & HZW 大结构有收益 \\
E10/E11 & E4/E8 路由融合 & 阈值或验证集路由。 & 阈值泛化不稳 \\
E12 & size-supervised dual slot & size 只进 loss 监督 slot2。 & HZW 强，MP 弱 \\
E13 & E12 + physics & 双槽位 + 物理约束。 & HZW 当前最好 \\
\bottomrule
\end{tabularx}
\end{frame}

\begin{frame}{E12 / E13 的设计逻辑}
\small
\begin{columns}[T,onlytextwidth]
\begin{column}{0.32\textwidth}
\begin{block}{结构}
以 ALIGNN 为骨架，保留 E4 的主超节点，再加入第二个 attentive global slot。
\end{block}
\end{column}
\begin{column}{0.33\textwidth}
\begin{block}{E12 损失}
\scriptsize
size 不进入前向 gate；训练时用 \texttt{num\_atoms} 构造软标签监督 slot2：
\[
\resizebox{0.95\linewidth}{!}{$
L = L_{\text{gap}} + \lambda \cdot \mathrm{BCE}\bigl(g,\sigma((N-80)/18)\bigr)
$}
\]
\end{block}
\end{column}
\begin{column}{0.32\textwidth}
\begin{block}{E13 物理约束}
在 E12 基础上加入 CBM/VBM 辅助监督和
\[
\text{gap} = \text{CBM} - \text{VBM}
\]
一致性约束，并用 mask 排除无效 CBM/VBM。
\end{block}
\end{column}
\end{columns}
\end{frame}

\section{HZW 结果}

\begin{frame}{HZW 阶段：模型增强确实有效}
\scriptsize
\centering
\resizebox{\textwidth}{!}{
\begin{tabular}{l l c c c l}
\toprule
实验 & 模型 & MAE$\downarrow$ & RMSE$\downarrow$ & $R^2\uparrow$ & 结论 \\
\midrule
E1 & CGCNN & 0.3022 & 0.5303 & 0.8847 & 基础 baseline \\
E2 & ALIGNN & 0.2314 & 0.4713 & 0.9089 & 最强 baseline \\
E4 & ALIGNN + VN & 0.2201 & 0.4746 & 0.9077 & 朴素超节点最稳 \\
E6 & ALIGNN + physics & 0.2212 & \textbf{0.4539} & \textbf{0.9155} & RMSE / $R^2$ 更好 \\
E8 & 2-slot attentive VN & 0.2285 & 0.4945 & 0.8997 & 整体一般，大结构有价值 \\
E12 & size-supervised dual slot & 0.2179 & 0.4712 & 0.9090 & 超过 E4 \\
E13 & E12 + physics & \textbf{0.2119} & 0.4616 & 0.9126 & HZW 当前最优 \\
\bottomrule
\end{tabular}
}

\vspace{1mm}
\small
HZW 上的核心事实：E4 证明“全局汇聚”有效；E12/E13 证明“带监督的第二全局槽位 + 物理约束”还能继续提升。
\end{frame}

\begin{frame}{HZW 分桶：为什么会想到双超节点}
\scriptsize
\centering
\setlength{\tabcolsep}{4pt}
\begin{tabular}{l c c c c l}
\toprule
原子数桶 & 样本数 & E4 & E7 & E8 & 观察 \\
\midrule
$\le 10$ & 2661 & \textbf{0.1248} & 0.2226 & 0.1272 & E4 最稳 \\
11--20 & 2232 & \textbf{0.2022} & 0.2912 & 0.2060 & E4 最稳 \\
21--40 & 2534 & \textbf{0.2424} & 0.3552 & 0.2548 & E4 最稳 \\
41--80 & 1642 & \textbf{0.3193} & 0.4419 & 0.3422 & E4 最稳 \\
81--160 & 447 & 0.3796 & 0.5190 & \textbf{0.3736} & E8 开始反超 \\
$>160$ & 35 & 0.2930 & 0.3462 & \textbf{0.2696} & E8 最好 \\
\bottomrule
\end{tabular}

\vspace{2mm}
\small
当时的合理推断：小中结构一个朴素超节点足够，大结构可能需要两个全局槽位承载不同信息。
\end{frame}

\section{MP 复现}

\begin{frame}{MP full 复现：规律开始失效}
\scriptsize
\centering
\begin{tabular}{l l c c c c}
\toprule
实验 & 模型 & MAE$\downarrow$ & RMSE$\downarrow$ & $R^2\uparrow$ & Best Epoch \\
\midrule
E2 & ALIGNN & 0.2397 & -- & -- & -- \\
E4 & ALIGNN + VN & \textbf{0.2302} & \textbf{0.5122} & \textbf{0.8854} & 50 \\
E7 & 1 attentive VN & 0.2390 & 0.5273 & 0.8785 & 48 \\
E8 & 2 attentive VN & 0.2418 & 0.5261 & 0.8791 & 43 \\
E12 & size-supervised & 0.2363 & 0.5187 & 0.8825 & 48 \\
\bottomrule
\end{tabular}

\vspace{2mm}
\small
\begin{itemize}
  \item MP full 上 E4 重新成为最优，E7/E8/E12 都没有超过 E4。
  \item 这说明 HZW 上的“size $\rightarrow$ 开第二槽位”假设不能直接迁移到 MP full。
  \item 但这还不能直接判定模型失败，因为目标标签来源已经发生变化。
\end{itemize}
\end{frame}

\begin{frame}{不是因为 MP 更小}
\scriptsize
\centering
\setlength{\tabcolsep}{4pt}
\begin{tabular}{l c c c c c c}
\toprule
数据集 & 样本数 & median N & mean N & $>80$ 占比 & $>160$ 占比 & 最大 N \\
\midrule
HZW all & 95,497 & 20 & 28.90 & 5.02\% & 0.35\% & 432 \\
MP full all & 154,875 & 20 & 31.39 & 7.76\% & 1.17\% & 444 \\
HZW test & 9,551 & 20 & 29.01 & 5.05\% & 0.37\% & 312 \\
MP test & 15,488 & 20 & 31.33 & 7.63\% & 1.16\% & 392 \\
\bottomrule
\end{tabular}

\vspace{2mm}
\small
\begin{itemize}
  \item MP 并不是更小，反而大结构占比更高。
  \item 因此“MP 小结构太多，所以双槽位无效”这个解释不成立。
  \item 更关键的问题是：\texttt{num\_atoms} 在 MP 标签下不是稳定的“需要第二全局槽位”的代理变量。
\end{itemize}
\end{frame}

\begin{frame}{MP-HZW matched clean 当前结果}
\scriptsize
\centering
\begin{tabular}{l l c c c c}
\toprule
实验 & 模型 & MAE$\downarrow$ & RMSE$\downarrow$ & $R^2\uparrow$ & Metal F1$\uparrow$ \\
\midrule
E2 & ALIGNN & 0.2534 & 0.5019 & 0.8973 & 0.8600 \\
E3 & M3GNet & 0.2966 & 0.5726 & 0.8663 & 0.8489 \\
E4 & ALIGNN + VN & \textbf{0.2446} & 0.5032 & 0.8968 & 0.8734 \\
E5 & Transformer & 0.2575 & 0.5179 & 0.8906 & 0.8670 \\
E7 & 1 attentive VN & 0.2478 & 0.5089 & 0.8944 & 0.8734 \\
E8 & 2 attentive VN & 0.2476 & 0.5111 & 0.8935 & \textbf{0.8757} \\
E10 & Val-routed E4/E8 & 0.2455 & 0.5044 & 0.8963 & 0.8730 \\
\bottomrule
\end{tabular}

\vspace{1mm}
\small
截至本地同步结果：E4 仍是最稳模型；E1/E12 clean 补跑结果尚未进入本地最终汇总。
\end{frame}

\begin{frame}{硬阈值路由为什么不稳}
\scriptsize
\centering
\setlength{\tabcolsep}{4pt}
\begin{tabular}{l c c c c l}
\toprule
原子数桶 & 样本数 & E4 & E8 & E10 & 现象 \\
\midrule
$\le10$ & 2525 & 0.1393 & \textbf{0.1378} & 0.1393 & E8 略好但 E10 没用上 \\
11--20 & 2086 & 0.2317 & \textbf{0.2276} & 0.2317 & E8 略好但 E10 没用上 \\
21--40 & 2300 & \textbf{0.2955} & 0.3008 & 0.2955 & E4 更好 \\
41--80 & 1603 & \textbf{0.3249} & 0.3359 & 0.3249 & E4 更好 \\
81--160 & 429 & \textbf{0.3364} & 0.3643 & 0.3620 & 阈值路由最吃亏 \\
$>160$ & 34 & 0.4598 & \textbf{0.3834} & \textbf{0.3834} & E8 明显更好 \\
\bottomrule
\end{tabular}

\vspace{2mm}
\small
验证集阈值 82 在 test 上不鲁棒；最优 test 阈值约 144，只保留极大结构给 E8。
\end{frame}

\section{数据质量诊断}

\begin{frame}{金属/非金属冲突：两个方向都存在}
\scriptsize
\centering
\setlength{\tabcolsep}{4pt}
\begin{tabularx}{\textwidth}{p{2.4cm} p{1.3cm} p{1.2cm} p{3.2cm} Y}
\toprule
冲突方向 & 数量 & 比例 & 强冲突例子 & 解释 \\
\midrule
MP metal-like，HZW nonmetal & 2,968 & 3.197\% & HZW gap $\ge 1$ eV: 1,130；$\ge 2$ eV: 447 & MP 认为近似金属，但 HZW 标签有明显带隙。 \\
MP nonmetal-like，HZW metal & 5,028 & 5.595\% & MP gap $\ge 1$ eV: 561；$\ge 2$ eV: 170 & MP 有正带隙，但 HZW 标签是 0。 \\
删除两类冲突后 & 84,830 & -- & train 67,865；val 8,481；test 8,484 & 得到 class-consistent matched 数据。 \\
\bottomrule
\end{tabularx}

\vspace{2mm}
\small
这不是简单噪声点，而是足以影响 metal/nonmetal 边界与回归损失的系统性标签冲突。
\end{frame}

\begin{frame}{同类样本仍有大数值差异}
\small
\begin{columns}[T,onlytextwidth]
\begin{column}{0.32\textwidth}
\begin{block}{数量}
删除金属/非金属冲突后，仍有 \textbf{1,317} 个 same-class 样本满足
\[
|\text{MP gap} - \text{HZW gap}| > 1\ \text{eV}.
\]
\end{block}
\end{column}
\begin{column}{0.32\textwidth}
\begin{block}{分布}
这些样本全部属于非金属类；其中 21--40 原子桶最多，为 \textbf{487} 个。
\end{block}
\end{column}
\begin{column}{0.32\textwidth}
\begin{alertblock}{含义}
即使金属/非金属类别一致，两个标签源在 gap 数值上也存在不可忽略偏差。
\end{alertblock}
\end{column}
\end{columns}
\end{frame}

\begin{frame}{当前最可信解释}
\small
\begin{itemize}
  \item HZW 上 E12/E13 有效，说明 dual-slot + size supervision + physics 的模型假设不是凭空无效。
  \item MP full 和 MP-HZW matched 上 E4 更稳，说明该假设对标签源和数据分布敏感。
  \item MP-HZW matched 不是“同一批结构、同一套标签”的纯复现实验；训练脚本使用的是 MP summary \texttt{band\_gap}，不是 HZW 原始 gap。
  \item 金属/非金属冲突和 $>1$ eV 数值差异会直接改变模型排序，因此当前负结果不能简单解释为模型结构失败。
\end{itemize}
\end{frame}

\section{结论}

\begin{frame}{总结}
\small
\begin{itemize}
  \item 完成了 CGCNN、ALIGNN、M3GNet 三类基线及多种 ALIGNN 增强模型的实验闭环。
  \item HZW 数据上：E4 证明朴素全局超节点有效，E12/E13 进一步取得最优结果。
  \item MP 数据上：复杂注意力超节点没有稳定超过 E4，说明原假设不能直接泛化。
  \item 数据诊断显示：MP 与 HZW 标签存在明显冲突，这是当前模型结论不一致的关键混杂因素。
\end{itemize}
\end{frame}

\begin{frame}[allowframebreaks]{参考文献}
\footnotesize
\begin{thebibliography}{99}
\bibitem{cgcnn} Tian Xie and Jeffrey C. Grossman. Crystal Graph Convolutional Neural Networks for an Accurate and Interpretable Prediction of Material Properties. \textit{Physical Review Letters}, 120(14):145301, 2018.
\bibitem{alignn} Kamal Choudhary and Brian DeCost. Atomistic Line Graph Neural Network for improved materials property predictions. \textit{npj Computational Materials}, 7:185, 2021.
\bibitem{m3gnet} Chi Chen and Shyue Ping Ong. A universal graph deep learning interatomic potential for the periodic table. \textit{Nature Computational Science}, 2:718--728, 2022.
\bibitem{mp} Anubhav Jain et al. Commentary: The Materials Project: A materials genome approach to accelerating materials innovation. \textit{APL Materials}, 1:011002, 2013.
\bibitem{matbench} Alexander Dunn et al. Benchmarking materials property prediction methods: the Matbench test set and Automatminer reference algorithm. \textit{npj Computational Materials}, 6:138, 2020.
\end{thebibliography}
\end{frame}

\end{document}
